\section*{Коротко о главном}
Слово \code{if} обозначает условие, а \code{if2} может быть именем переменной. Сканер обязан различать их при любом устройстве программы. Мы задали четырём сканерам одни правила, проверили выдачу по каждому байтовому интервалу и измерили полный проход по 21 тексту. За единицу сравнения взяли время прохода, делённое на число байтов входа.

\begin{enumerate}
\item \textbf{На учебном псевдокоде быстрее всего оказался re2c.}
Его медиана составила 2,478~нс/байт. Ближайший результат у собственного варианта с ветвлениями, 3,432~нс/байт. На смеси из повторений строки с условным оператором наименьшее время при всех трёх объёмах показал табличный сканер.
\item \textbf{Длина лексемы меняет результат даже при одинаковом объёме файла.}
На идентификаторах из двух букв таблица выигрывала у ветвлений во всех 16 раундах. На идентификаторах из 32 и 128 букв во всех раундах выигрывали ветвления. Для длины 128 медиана парного отношения времён достигла 3,245.
\item \textbf{Близкие результаты требуют оговорки.} На повторениях \code{abcde2} медианы \code{table} и \code{re2c} равны 1,551 и 1,555~нс/байт. Межквартильные интервалы перекрываются. На этом входе данные не дают основания предпочесть одну из этих программ по скорости.
\end{enumerate}
В конце даны три рекомендации по выбору и сравнению сканеров.

\subhead{Что получает сканер}
Первая строка корпуса взята из записи функции модуля числа:
\begin{center}
\code{abs(x) := if x>0 then x else -x end if}
\end{center}
Сканер выдаёт вид и границы каждой лексемы. Получатель может отличить имя \code{abs} от ключевого слова \code{if}, найти присваивание и извлечь исходное написание числа. Для этой задачи не требуется проверять синтаксис условного оператора или вычислять значение функции \sourcega{3.2.1}{1/4}.

\begin{center}
\captionof{table}{Лексика выбранного псевдокода. В первых двух строках приведены регулярные выражения, в остальных~--- полный перечень написаний.}
\small\renewcommand{\arraystretch}{1.18}
\begin{tabularx}{\linewidth}{@{}lX@{}}\toprule
Идентификатор & \code{[A-Za-z][A-Za-z0-9]*}\\
Десятичное число & \code{[0-9]+}\\
Ключевые слова & \code{if then else end while do for from to return case}\\
Знаки и разделители & \code{:= \ = \ != \ < \ <= \ > \ >= \ + \ - \ * \ / \ ( \ ) \ [ \ ] \ , \ ; \ :}\\ \bottomrule
\end{tabularx}
\end{center}
